\appendix
\section{Anexos}

\subsection{Instrumento de validación}

El instrumento completo se entrega en el archivo
\texttt{encuesta\_validacion.md}. Está diseñado como un formulario con rutas para
responsables y miembros adultos. La meta es obtener 10 respuestas válidas por
ruta. Los resultados y gráficos se incorporarán después de su aplicación.

\subsection{Evidencia técnica disponible}

\begin{table}[H]
\caption{Inventario de evidencia técnica}
\begin{tabularx}{\textwidth}{>{\bfseries}p{4cm}X p{3cm}}
\toprule
Evidencia & Descripción & Estado \\
\midrule
Repositorio & Código frontend, backend, infraestructura y documentación. & Disponible \\
Pruebas backend & 18 archivos y 148 casos identificados. & Disponible \\
Pruebas frontend & 26 archivos de pruebas unitarias o de componentes. & Disponible \\
Pruebas E2E & Dos especificaciones Playwright. & Disponible \\
Frontend público & Despliegue demostrativo en Vercel. & Disponible \\
Backend y API & Despliegue y documentación OpenAPI. & Disponible \\
Capturas anonimizadas & Pantallas seleccionadas para el escrito o cartel. & Pendiente \\
Video offline & Recorrido de respaldo para la exposición. & Pendiente \\
\bottomrule
\end{tabularx}
\caption*{\textit{Nota}. Elaboración propia a partir del repositorio al 1 de julio de 2026 \parencite{gymhub_repo}.}
\end{table}

\subsection{Pendientes controlados}

\begin{itemize}
  \item Confirmar CORVEC en la portada.
  \item Consultar al CTR la aplicación de ExpoTEC-5 y ExpoTEC-6.
  \item Aplicar la encuesta y documentar fecha, muestra, resultados y limitaciones.
  \item Obtener una cotización de dominio y completar la valoración de horas.
  \item Incorporar capturas anonimizadas y video offline.
  \item Actualizar tarifas tecnológicas antes de entregar o vender.
\end{itemize}

\subsection{Referencias de despliegue demostrativo}

\begin{itemize}
  \item Frontend: \url{https://proyectoappgym-frontend.vercel.app}
  \item Backend: \url{https://proyectoappgym-backend.vercel.app}
  \item Documentación API: \url{https://proyectoappgym-backend.vercel.app/api/docs/}
\end{itemize}

No se incluyen credenciales en el documento público. La demostración utilizará
cuentas preparadas por el equipo y datos ficticios.
